\chapter{Graph Colouring}

\section{Vertex colourings}

\begin{notedefinition}
Let $G=(V,E)$ be a graph.  A \emph{proper vertex colouring} is a map
$c\colon V\to S$ such that $c(u)\ne c(v)$ whenever $uv\in E$.  If
$|S|=k$, it is a \emph{$k$-colouring}.  The least such $k$ is the
\emph{chromatic number} $\chi(G)$.
\end{notedefinition}

\begin{notetheorem}[Four-Colour Theorem]
Every planar graph is $4$-colourable.
\end{notetheorem}

\begin{notetheorem}[Five-Colour Theorem]
Every planar graph is $5$-colourable.
\end{notetheorem}
\begin{proof}
Suppose a counterexample $G$ has as few vertices as possible, and fix a plane
embedding.  By Corollary~7.5, choose a vertex $v$ with $\deg(v)\le5$.  By
minimality, $G-v$ has a $5$-colouring.  If at most four colours occur on
$N(v)$, an unused colour extends the colouring to $v$.

It remains to consider $\deg(v)=5$ with all five colours appearing.  List the
neighbours $v_1,v_2,v_3,v_4,v_5$ in cyclic order around $v$, with
$c(v_i)=i$.  Let $H_{ij}$ be the subgraph induced by vertices of colours $i$
and $j$.

If $v_1$ and $v_3$ lie in different components of $H_{13}$, interchange
colours $1$ and $3$ in the component containing $v_1$.  Colour $1$ is then
absent from $N(v)$ and may be assigned to $v$.

Otherwise there is a $1$--$3$ Kempe path from $v_1$ to $v_3$.  Together with
$v v_1$ and $v v_3$, it forms a closed curve separating $v_2$ from $v_4$.
Hence no $2$--$4$ Kempe path joins $v_2$ to $v_4$.  Interchanging colours
$2$ and $4$ in the component containing $v_2$ frees colour $2$ for $v$.
Both cases contradict the choice of $G$.
\end{proof}
\begin{notetheorem}[greedy colouring]
For every graph $G$,
\[
  \chi(G)\le\Delta(G)+1.
\]
\end{notetheorem}
\begin{proof}
Order the vertices arbitrarily and colour them one at a time.  When a vertex
is reached, at most $\Delta(G)$ colours are forbidden by already coloured
neighbours, so one of $\Delta(G)+1$ colours remains available.
\end{proof}

\begin{notelemma}[Brooks ordering lemma]
Let $G$ be a connected graph with maximum degree $\Delta\ge3$.  If $G$ is
neither complete nor an odd cycle, then its vertices can be ordered so that
\begin{itemize}
  \item two nonadjacent vertices $a,b$ are coloured first with the same
        colour;
  \item every later vertex except the final one has a neighbour still
        uncoloured when it is coloured;
  \item the final vertex is adjacent to both $a$ and $b$.
\end{itemize}
\end{notelemma}
\begin{proof}[Proof sketch]
For a $2$-connected graph that is neither complete nor a cycle, one can choose
vertices $a,b$ at distance two such that $G-\{a,b\}$ is connected; let $v$ be
a common neighbour.  Take a spanning tree of $G-\{a,b\}$ rooted at $v$ and
order its vertices from leaves toward the root.  This gives the required
ordering.  If $G$ has a cut-vertex, apply the same construction successively
to end blocks, permuting colour names at the shared cut-vertex.  Complete end
blocks cannot have order $\Delta+1$, because their cut-vertex has an
additional neighbour outside the block.  This standard block reduction leaves
only complete graphs and odd cycles as the exceptional connected graphs.
\end{proof}

\begin{notetheorem}[Brooks's Theorem]
If a connected graph $G$ is neither complete nor an odd cycle, then
\[
  \chi(G)\le\Delta(G).
\]
\end{notetheorem}
\begin{proof}
The assertion is immediate for $\Delta(G)\le2$.  For $\Delta\ge3$, use the
ordering from the lemma.  Give the nonadjacent vertices $a,b$ the same colour.
Colour the remaining vertices greedily in the prescribed order using
$\Delta$ colours.  Every nonfinal vertex has an uncoloured neighbour, so at
most $\Delta-1$ colours are forbidden.  At the final vertex, the two
neighbours $a,b$ use the same colour, so again at most $\Delta-1$ distinct
colours appear.  Thus the colouring succeeds with $\Delta$ colours.
\end{proof}
\section{Edge colourings}

\begin{notedefinition}
A \emph{proper edge colouring} is a map $c\colon E(G)\to S$ such that edges
sharing a vertex receive different colours.  The least possible number of
colours is the \emph{chromatic index} $\chi'(G)$.
\end{notedefinition}

\begin{noteremark}
Every colour class in a proper edge colouring is a matching, and
\[
  \chi'(G)\ge\Delta(G).
\]
A graph is $2$-edge-colourable exactly when each nontrivial component is a
path or an even cycle.
\end{noteremark}

\begin{notelemma}[Vizing fan extension lemma]
Let $G$ be simple, let $uv$ be an uncoloured edge, and suppose all other edges
have been properly coloured with $\Delta(G)+1$ colours.  By rotating colours
on a maximal fan centred at $u$ and, if necessary, interchanging two colours
along a maximal alternating Kempe chain, the colouring can be extended to
$uv$.
\end{notelemma}
\begin{proof}[Proof sketch]
Build a maximal fan
$v=v_0,v_1,\dots,v_k$ at $u$, where the colour on $uv_i$ is missing at an
earlier fan vertex.  Let $\alpha$ be missing at $u$ and let $\beta$ be missing
at $v_k$.  Consider the maximal $\alpha$--$\beta$ alternating chain through
$v_k$.  If it avoids $u$, swap $\alpha$ and $\beta$ on that chain and rotate
the fan, making one colour free at both ends of the uncoloured edge.  If it
reaches $u$, the defining missing colours of the fan locate an earlier fan
vertex at which a second Kempe interchange followed by a shorter fan rotation
has the same effect.  Maximality of the fan guarantees that one of these two
operations applies.  All operations preserve properness because they merely
swap colours on a two-coloured component or move a missing colour along the
fan.
\end{proof}

\begin{notetheorem}[Vizing's Theorem]
For every simple graph $G$,
\[
  \Delta(G)\le\chi'(G)\le\Delta(G)+1.
\]
\end{notetheorem}
\begin{proof}
The lower bound is immediate at a vertex of maximum degree.  For the upper
bound, induct on $|E(G)|$.  Remove an edge $uv$ and colour $G-uv$ with
$\Delta(G)+1$ colours by induction.  The Vizing fan extension lemma extends
this colouring to $uv$.
\end{proof}

\begin{noteremark}
Simple graphs with $\chi'(G)=\Delta(G)$ are called \emph{class 1}; those with
$\chi'(G)=\Delta(G)+1$ are \emph{class 2}.  The theorem is false in this form
for multigraphs, where more than $\Delta+1$ colours may be needed.
\end{noteremark}
\begin{notetheorem}[K\H{o}nig's Line-Colouring Theorem]
If $G$ is bipartite, then
\[
  \chi'(G)=\Delta(G).
\]
\end{notetheorem}
\begin{proof}
Let $\Delta=\Delta(G)$.  Add dummy vertices so that the two bipartition
classes have the same size, and then add edges between deficient vertices
until every vertex has degree $\Delta$.  Parallel edges may be used in this
auxiliary bipartite multigraph $R$; the total degree deficit is the same on
both sides, so the process terminates.

Every regular bipartite multigraph has a perfect matching by Hall's theorem.
Remove one perfect matching from $R$; the remaining graph is
$(\Delta-1)$-regular.  Repeating decomposes $E(R)$ into $\Delta$ perfect
matchings.  Give each matching its own colour and restrict the colouring to
$G$.  Thus $\chi'(G)\le\Delta$, while the reverse inequality always holds.
\end{proof}
\section{Exercises and extremal examples}

\begin{noteexercise}
Suppose every two odd cycles of a graph $G$ have a common vertex.  Then
$G$ is $5$-colourable.
\end{noteexercise}
\begin{proof}
If $G$ has no odd cycle, it is bipartite and hence $2$-colourable.  Otherwise
fix an odd cycle $C$.  Every odd cycle meets $C$, so $G-V(C)$ has no odd cycle
and is bipartite.  Colour $G-V(C)$ with two colours and colour $C$ with three
new colours.  Since the two colour sets are disjoint, this is a proper
$5$-colouring of $G$.
\end{proof}
\begin{noteexercise}
For the complete graph $K_n$,
\[
  \chi'(K_n)=
  \begin{cases}
    n-1,& n\text{ even},\\
    n,& n\text{ odd}.
  \end{cases}
\]
\end{noteexercise}
\begin{proof}
The lower bound $\chi'(K_n)\ge n-1$ is the degree bound.  If $n$ is odd, each
matching contains at most $(n-1)/2$ edges, while
$|E(K_n)|=n(n-1)/2$; therefore at least $n$ colours are needed.

For odd $n$, label the vertices by $\mathbb Z_n$ and colour the edge $ij$ by
$i+j\pmod n$.  At a fixed vertex $i$, the values $i+j$ are distinct as
$j$ varies, so this is a proper $n$-edge-colouring.

For even $n=2m$, label the vertices by
$\{\infty\}\cup\mathbb Z_{2m-1}$.  For each
$c\in\mathbb Z_{2m-1}$, let
\[
  M_c
  :=\{\infty c\}
    \cup\bigl\{(c+i)(c-i):1\le i\le m-1\bigr\},
\]
with arithmetic modulo $2m-1$.  Each $M_c$ is a perfect matching, and the
$2m-1$ matchings partition $E(K_{2m})$.  Assigning one colour to each $M_c$
gives an $(n-1)$-edge-colouring.
\end{proof}
\begin{noteexercise}
Let $G$ have $n$ vertices and let $\overline G$ be its complement.  Then
\[
  \chi(G)\,\chi(\overline G)\ge n.
\]
\end{noteexercise}
\begin{proof}
Choose proper colourings
$p\colon V(G)\to[\chi(G)]$ and
$q\colon V(G)\to[\chi(\overline G)]$.  Assign each vertex the ordered pair
$(p(v),q(v))$.  Distinct vertices cannot receive the same pair: if they did,
they would be nonadjacent in $G$ because their $p$-colours agree and
nonadjacent in $\overline G$ because their $q$-colours agree, which is
impossible.  The assignment is injective, so the number of available pairs is
at least $n$.
\end{proof}
